\documentclass[12pt,a4paper,twoside]{article}
\usepackage{graphicx} % Required for inserting images
\usepackage{fancyhdr}
\usepackage{indentfirst}
\usepackage{amsmath}   % Entornos matemáticos avanzados
\usepackage{amsfonts} % Conjuntos como \mathbb{R}
\usepackage{amssymb}  % Símbolos extra
\usepackage{float}
\usepackage{listings}
\usepackage{subcaption}
\usepackage[spanish]{babel}
\usepackage{dirtree}
\usepackage{etoolbox}
\usepackage{pdfpages}
\usepackage[hidelinks]{hyperref}
\usepackage[babel]{csquotes} % muy recomendado con biblatex
\usepackage[  backend=biber,
  style=numeric,
  sorting=nyt]{biblatex}
\usepackage{makecell}
\DeclareLanguageMapping{spanish}{spanish-apa}
\addbibresource{bibliografia.bib}

\usepackage[left=2.5cm, right=2.5cm, top=3cm, bottom=3cm]{geometry}
% ------------------------------
% Estilo de títulos (bloques visuales)
% ------------------------------
\usepackage{titlesec}
\titleformat{\chapter}[hang]
  {\Huge\bfseries}{\thechapter}{2em}{}[\titlerule] % Línea bajo capítulos
\titleformat{\section}[hang]
  {\Large\bfseries}{\thesection}{1em}{}[\titlerule] % Línea bajo secciones
\titleformat{\subsection}[hang]
  {\large\bfseries}{\thesubsection}{1em}{} % sin línea
\titleformat{\subsubsection}[hang]
  {\normalsize\bfseries}{\thesubsubsection}{1em}{} % sin línea

% ------------------------------
% Encabezados y pies
% ------------------------------
\usepackage{fancyhdr}
\pagestyle{fancy}
\fancyhf{} % limpia
\fancyhead[L]{\leftmark} % Nombre de capítulo/sección a la izquierda
\fancyhead[R]{\thepage}  % Número de página a la derecha

\title{Deep Learning\\ Sistema de Mantenimiento Predictivo}
\author{Libe Belasko, Javier Solana, Miguel Iriso}
\date{November 2025}

\begin{document}
\includepdf[
  pages={1-2},
  fitpaper=true
]{portada.pdf}
% ===== PÁGINAS PRELIMINARES (sin números o con números romanos) =====
%\mbox{}
\newpage
\pagestyle{empty}
\section*{Resumen}

\newpage
\pagenumbering{roman}
\pagestyle{plain}
%Índice
\tableofcontents
\newpage
\listoffigures  % Índice de figuras
\newpage
\listoftables   % Índice de tablas
\newpage
%Primera seccion
\pagenumbering{arabic} % Números arábigos
% Configurar el estilo de numeración
\pagestyle{plain}

\pretocmd{\section}{\cleardoublepage}{}{}
%\pagestyle{fancy}
%\fancyhf{} % Limpiar encabezados y pies
%\fancyfoot[C]{\thepage} % Número de página centrado

\section{Introducción}
El rápido avance de la inteligencia artificial generativa ha transformado profundamente la manera
en que se crea contenido digital. Herramientas que hace apenas unos años estaban reservadas a
estudios de producción especializados son hoy accesibles de forma gratuita en internet, permitiendo
a cualquier persona generar imágenes fotorrealistas prácticamente indistinguibles de una fotografía
real. Este fenómeno ha dado lugar a una categoría creciente de medios sintéticos comúnmente
denominados imágenes generadas por IA, producidas mediante arquitecturas como las redes
generativas adversariales (GANs), los modelos de difusión o los transformadores visuales de gran
escala.

Si bien estas tecnologías tienen un valor creativo y práctico innegable, su uso indebido plantea
riesgos serios que van mucho más allá del ámbito del entretenimiento. La falsificación de documentos
de identidad, la creación de retratos faciales fraudulentos para ataques de ingeniería social, o
la fabricación de pruebas visuales en contextos legales e institucionales son solo algunos de los
escenarios en los que la incapacidad de distinguir lo real de lo sintético puede acarrear
consecuencias directas tanto para personas como para organizaciones. El problema se agrava por la
velocidad a la que mejora la calidad de generación: modelos publicados con apenas meses de
diferencia pueden dejar obsoletos los sistemas de detección anteriores.

Los métodos de detección más avanzados han demostrado una precisión prometedora en entornos
controlados; sin embargo, en su mayoría funcionan como clasificadores opacos. Un modelo que
devuelve una puntuación de confianza sin ninguna justificación adicional ofrece una utilidad limitada en situaciones de alto riesgo. Un
analista forense, un periodista o un profesional del ámbito legal no puede actuar únicamente sobre
un número: necesita entender \textit{por qué} el sistema ha llegado a esa conclusión. ¿Qué regiones
de la imagen han activado la decisión? ¿Se trata de una textura anómala en la piel, una
inconsistencia en la iluminación o artefactos en los bordes del cabello? Sin esta información, el
sistema de detección no puede ser auditado, cuestionado ni, en última instancia, confiable.

Esta brecha entre rendimiento predictivo e interpretabilidad es la que motiva el enfoque central
de este trabajo. En lugar de limitarnos a determinar \textit{si} una imagen ha sido generada por
IA, nuestro objetivo es estudiar aproximaciones que también sean capaces de explicar \textit{dónde}
y, en cierta medida, \textit{por qué}. Mediante la incorporación de técnicas de explicabilidad
---como mapas de atribución basados en gradientes o métodos de oclusión--- sobre modelos de
detección existentes, buscamos producir resultados que sean a la vez precisos e interpretables.
Este enfoque se enmarca dentro del movimiento hacia una IA confiable (\textit{Trustworthy AI}),
que reconoce la transparencia y la supervisión humana como propiedades esenciales de cualquier
sistema desplegado en contextos reales sensibles.

El resto del artículo presenta una revisión de la literatura relevante, una comparativa de los
enfoques de detección y explicabilidad seleccionados, y una evaluación empírica diseñada para
medir tanto la precisión de los modelos como la calidad de las explicaciones generadas.


\section{Estado del Arte}

\subsection{Detección de imágenes generadas por IA}

La detección de imágenes sintéticas ha evolucionado de forma acelerada entre 2022 y
2025, pasando de clasificadores superficiales basados en características artesanales
a arquitecturas profundas que explotan representaciones de modelos fundacionales. A
grandes rasgos, los enfoques actuales se agrupan en tres líneas: detectores
supervisados con CNN o ViT, métodos basados en señales de reconstrucción o
gradiente, y detectores que aprovechan espacios de características de modelos
preentrenados sin necesidad de entrenamiento adicional.

\textbf{Detectores supervisados basados en CNN.} Los backbones convolucionales
(EfficientNet, Xception, ConvNeXt) preentrenados en ImageNet y ajustados sobre
datos reales frente a sintéticos constituyen la línea de trabajo más madura,
alcanzando AUC $\geq 0.99$ en benchmarks de distribución cerrada como
FaceForensics++. Sin embargo, su rendimiento cae de forma notable ante generadores
no vistos durante el entrenamiento. Wang et al.~\cite{wang2022dual} demuestran que
añadir una rama en el dominio frecuencial (FFT, filtros SRM) mejora la robustez
frente a compresión y redimensionado, capturando los artefactos periódicos que
dejan los pasos de \textit{up-sampling} de las GAN.

\textbf{Representaciones de gradiente y error de reconstrucción.} Tan et
al.~\cite{tan2023lgrad} proponen \textbf{LGrad} (CVPR 2023), uno de los trabajos
más citados del período con más de 190 referencias en Semantic Scholar. En lugar de
operar directamente sobre píxeles RGB, calculan mapas de gradiente de la imagen
respecto a una CNN preentrenada; estos gradientes capturan artefactos de generación
universales que se transfieren entre arquitecturas GAN distintas y conjuntos de
datos diferentes. En una dirección complementaria, Wang et al.~\cite{wang2023dire}
introducen \textbf{DIRE} (ICCV 2023), un detector específico para imágenes de
difusión. La intuición es que las imágenes generadas por un modelo de difusión se
reconstruyen con bajo error cuando se pasan de nuevo por ese mismo modelo, mientras
que las imágenes reales presentan un error de reconstrucción alto. DIRE ofrece buena
generalización a modelos de difusión no vistos y viene acompañado del benchmark
\textit{DiffusionForensics}.

\textbf{Detectores basados en modelos preentrenados.} Ojha et al.~\cite{ojha2023univfd}
presentan \textbf{UnivFD} (CVPR 2023), argumentando que los clasificadores
supervisados estándar sobreajustan a los generadores vistos durante el
entrenamiento. Su propuesta es realizar la clasificación real-vs-falso mediante
vecinos más cercanos en el espacio de características de CLIP, sin ajustar el
backbone. Este planteamiento mejora la precisión media en aproximadamente 26 puntos
porcentuales frente a detectores supervisados cuando se evalúa sobre 19 modelos
generativos no vistos, incluyendo GAN y difusión. Ampliando esta idea, Sha et
al.~\cite{sha2023raising} muestran que un ViT-B/16 de CLIP con una cabeza lineal
ligera alcanza AUC $\geq 0.99$ en benchmarks GAN estándar y mantiene entre 0.92 y
0.96 ante generadores desconocidos.

Más recientemente, \textbf{RIGID}~\cite{rigid2024} propone un detector sin
entrenamiento (\textit{training-free}) basado en la observación de que las imágenes
reales son más robustas que las sintéticas ante pequeñas perturbaciones en el
espacio de representaciones de modelos fundacionales como \textbf{DINOv2} o CLIP.
La caída de similitud entre el embedding original y el perturbado actúa como señal
de detección, sin requerir datos de entrenamiento adicionales ni ajuste fino. RIGID
demuestra resultados competitivos en GenImage y otros benchmarks modernos,
convirtiendo a DINOv2 en una pieza relevante del ecosistema de detección actual.

\textbf{Benchmarks de referencia.} Zhu et al.~\cite{zhu2023genimage} introducen
\textbf{GenImage} (NeurIPS 2023), un benchmark de más de un millón de pares
real-sintético producidos por más de siete familias generativas (StyleGAN,
Stable Diffusion, DALL-E, entre otras). Define dos tareas principales: clasificación
cruzada entre generadores y clasificación con degradaciones (baja resolución,
compresión). Se ha convertido en el benchmark de referencia para evaluar
generalización, y sus resultados muestran que incluso los mejores modelos actuales
caen por debajo de AUC 0.90 cuando el generador de test es completamente desconocido.

\subsection{Explicabilidad aplicada a la detección}

Dado que los detectores de alto rendimiento operan en gran medida como cajas negras,
ha surgido una línea específica que aplica técnicas de IA explicable (XAI) para
justificar sus decisiones, especialmente relevante en contextos forenses y legales.

Los \textbf{métodos basados en gradiente}, y en particular Grad-CAM y sus variantes,
son los más utilizados por su bajo coste y fácil integración en CNN. En detección de
deepfakes se aplican sobre la última capa convolucional para resaltar las regiones
faciales que más han influido en la predicción. Silva et al.~\cite{silva2022}
proponen un ensemble jerárquico donde los mapas de atención localizan artefactos en
zonas de boca, ojos y bordes faciales en FaceForensics++. Sistemas como el de Arjun
et al.~\cite{arjun2023xception} combinan XceptionNet con Grad-CAM y MediaPipe Face
Mesh para traducir mapas de calor en descriptores anatómicos de alto nivel,
facilitando la interpretación por parte de analistas humanos. Trabajos más recientes
proponen ensembles de MobileNetV2 y EfficientNetV2B0 donde Grad-CAM actúa
simultáneamente como mecanismo de interpretabilidad y herramienta de
depuración~\cite{gradcam_mobile}.

Los \textbf{métodos basados en perturbaciones} (LIME, SHAP, RISE) ofrecen
estimaciones de relevancia local modificando partes de la entrada y observando el
impacto en la predicción. Aunque más costosos que Grad-CAM, permiten validar la
fidelidad causal de las explicaciones. El marco de evaluación cuantitativa de
\citeauthor{quant2024}~\cite{quant2024} compara cinco métodos aplicados a un
detector entrenado en FaceForensics++, concluyendo que LIME obtiene mayor fidelidad
causal medida por las métricas de \textit{insertion/deletion}, mientras que los
métodos gradiente producen visualizaciones perceptualmente más intuitivas. El
sistema \textbf{DeepExplain}~\cite{deepexplain2024} combina ambas familias en un
pipeline unificado: las saliencias de Grad-CAM alimentan el cálculo de valores SHAP
a nivel de región, reportando además mejoras en AUROC respecto a baselines sin
módulo explicativo.

Un resultado transversal de esta literatura es que las métricas cualitativas y
cuantitativas de las explicaciones no siempre están alineadas: mapas de calor
visualmente convincentes pueden no reflejar fielmente el proceso interno de
decisión, lo que resulta crítico en contextos forenses donde las explicaciones
pueden tener implicaciones legales~\cite{gradcam_medical}. Esto motiva que los
trabajos más recientes apunten hacia sistemas que validen las explicaciones mediante
métricas de fidelidad y localización, y que evalúen explícitamente su robustez
adversarial.

\section{Metodología}

\subsection{Visión general}

El objetivo de este trabajo es comparar tres aproximaciones distintas para la
detección de imágenes generadas por IA, evaluando no solo su rendimiento predictivo
sino también la calidad e interpretabilidad de las explicaciones que producen. Para
ello se sigue un protocolo común: los tres modelos se evalúan sobre el mismo
conjunto de test, se aplica una técnica de explicabilidad sobre cada uno, y los
resultados se comparan mediante métricas cuantitativas y análisis visual cualitativo.

\subsection{Dataset}

Se utiliza \textbf{CIFAKE}~\cite{cifake}, un conjunto de datos público disponible en
Kaggle que contiene 60.000 imágenes divididas en dos clases: imágenes reales
procedentes de CIFAR-10 e imágenes sintéticas generadas con Stable Diffusion
condicionadas a las mismas categorías semánticas. La distribución es equilibrada
(50\% real, 50\% sintético), lo que evita sesgos de clase en las métricas de
evaluación.

Para mantener los experimentos dentro de los límites de memoria y tiempo de
ejecución de Google Colab, se trabaja con un subconjunto de \textbf{10.000 imágenes}
(8.000 para entrenamiento y validación, 2.000 para test). El conjunto de test es
idéntico para los tres modelos, garantizando comparabilidad directa de los
resultados. Todas las imágenes se redimensionan a $224 \times 224$ píxeles y se
normalizan con la media y desviación estándar de ImageNet.

\subsection{Modelos de detección}

\subsubsection{EfficientNet-B4 supervisado}

Como baseline convolucional se emplea \textbf{EfficientNet-B4}~\cite{tan2019efficientnet},
preentrenado en ImageNet. Se sustituye la cabeza de clasificación original por una
capa lineal binaria y se realiza \textit{fine-tuning} completo durante 10 épocas con
optimizador Adam (lr $= 1\text{e-}4$, weight decay $= 1\text{e-}5$) y
\textit{early stopping} sobre la pérdida de validación. Este modelo representa la
familia de detectores supervisados que domina los benchmarks de distribución cerrada.

\subsubsection{UnivFD — clasificación mediante CLIP}

El segundo modelo sigue el planteamiento de Ojha et al.~\cite{ojha2023univfd}: se
congela el encoder visual \textbf{ViT-B/16 de CLIP} y se entrena únicamente una
cabeza lineal binaria sobre sus representaciones. El backbone no se ajusta, lo que
evita el sobreajuste al generador específico presente en el dataset de entrenamiento
y favorece la generalización. El entrenamiento de la cabeza lineal converge en menos
de 5 épocas con el mismo optimizador y tasa de aprendizaje que EfficientNet.

\subsubsection{RIGID — detección sin entrenamiento con DINOv2}

El tercer modelo implementa el enfoque de \textbf{RIGID}~\cite{rigid2024}: se
emplea el encoder \textbf{DINOv2 ViT-B/14}~\cite{oquab2023dinov2}, entrenado de
forma autosupervisada, sin ningún ajuste fino adicional. La señal de detección se
construye perturbando los embeddings con ruido gaussiano de pequeña magnitud
($\sigma = 0.01$) y midiendo la caída en similitud coseno entre el embedding
original y el perturbado; las imágenes sintéticas presentan una caída
significativamente mayor que las reales. La clasificación final se realiza mediante
un umbral sobre esta diferencia, optimizado sobre el conjunto de validación. Este
modelo no requiere ninguna fase de entrenamiento.

\subsection{Métodos de explicabilidad}

\subsubsection{Grad-CAM sobre EfficientNet}

Se aplica \textbf{Grad-CAM}~\cite{selvaraju2017gradcam} sobre la última capa
convolucional de EfficientNet-B4, generando un mapa de calor que resalta las
regiones de la imagen cuya activación ha contribuido más a la predicción de clase
sintética. Los mapas se superponen sobre la imagen original con una escala de color
continua para facilitar la inspección visual.

\subsubsection{Mapas de atención de DINOv2}

Una ventaja intrínseca de los modelos entrenados con DINO es que sus cabezas de
\textit{self-attention} producen mapas de atención interpretables sin necesidad de
post-proceso adicional~\cite{caron2021dino}. Se extraen los mapas de atención del
último bloque transformer de DINOv2 y se promedian entre cabezas para obtener una
sola visualización por imagen.

\subsubsection{LIME sobre UnivFD}

Para el modelo CLIP-based se aplica \textbf{LIME}~\cite{ribeiro2016lime}
(\textit{Local Interpretable Model-agnostic Explanations}), que segmenta la imagen
en super-píxeles mediante SLIC y estima la contribución de cada región a la
predicción entrenando un modelo lineal local sobre perturbaciones de la entrada. Se
generan 1.000 muestras por imagen y se retienen los 10 super-píxeles más relevantes.

\subsection{Métricas de evaluación}

\textbf{Rendimiento predictivo.} Se reportan \textit{accuracy}, AUC-ROC y F1-score
para los tres modelos sobre el conjunto de test común.

\textbf{Calidad de las explicaciones.} Se evalúa mediante la métrica de
\textit{insertion/deletion}~\cite{petsiuk2018rise}: se insertan progresivamente las
regiones más relevantes según cada explicador y se mide cómo evoluciona la confianza
del modelo; una subida rápida indica explicaciones de alta fidelidad. Adicionalmente,
se realiza un análisis visual cualitativo sobre una muestra de 20 imágenes
seleccionadas de forma estratificada.



\end{document}

